% ch8-jordan.tex — Chapter 8: Jordan Normal Form
% Nilpotent endomorphisms, generalized eigenspaces, Jordan blocks,
% the Jordan normal form theorem, matrix exponential, applications.

\chapter{Jordan Normal Form}\label{ch:jordan}

% ============================================================
% Motivating introduction
% ============================================================

In 
ef{ch:eigenvalues} we saw that a matrix is diagonalizable if and
only if the algebraic and geometric multiplicities of every eigenvalue
coincide.  When they do, one obtains the simplest possible matrix
representation --- a diagonal matrix.  But what happens when
diagonalization \emph{fails}?

Consider the matrix
$A = \begin{psmallmatrix} 2 & 1 \\ 0 & 2 \end{psmallmatrix}$.
Its only eigenvalue is $\lambda = 2$ with algebraic multiplicity~$2$,
yet the eigenspace $E_2(A) = \Span\!\bigl(\begin{psmallmatrix} 1 \\ 0
\end{psmallmatrix}\bigr)$ has dimension~$1$.  There is no basis
of~$\R^2$ consisting of eigenvectors of~$A$, so $A$ is not
diagonalizable.  Nevertheless, $A$ is already in a rather pleasant
form: it is ``almost diagonal,'' differing from $\diag(2,2)$ by a
single off-diagonal entry.  Is this an accident, or can every matrix be
brought into a similarly nice shape?

The answer is provided by the \emph{Jordan normal form}\index{Jordan
normal form}, one of the most powerful results in linear algebra.  It
asserts that over an algebraically closed field (such as~$\C$), every
square matrix is similar to an essentially unique matrix built from
simple building blocks called \emph{Jordan blocks}.  This canonical form
completely determines the matrix up to similarity and provides a
systematic way to handle all the phenomena that arise when
diagonalization fails.

The Jordan form has far-reaching applications: it yields explicit
formulas for matrix powers and matrix exponentials, it explains the
structure of solutions to linear systems of differential equations with
non-diagonalizable coefficient matrices, and it clarifies the
relationship between the characteristic and minimal polynomials.

Throughout this chapter, unless stated otherwise, $\K$ denotes an
algebraically closed field (typically $\K = \C$), $E$ is a
finite-dimensional $\K$-vector space with $\dim E = n$, and
$f \in \End(E)$.  We freely use results from 
ef{ch:eigenvalues},
especially the notions of eigenvalue, eigenspace, characteristic
polynomial, and minimal polynomial.


% ============================================================
\section{Nilpotent endomorphisms}\label{sec:nilpotent}
% ============================================================

The key to understanding non-diagonalizable maps is to understand a
special class of endomorphisms that are ``as far from diagonalizable as
possible'' (aside from the zero map): the \emph{nilpotent}
endomorphisms.

\begin{definition}{Nilpotent endomorphism}{def:nilpotent}
\index{nilpotent!endomorphism}\index{endomorphism!nilpotent}
An endomorphism $f \in \End(E)$ is called \emph{nilpotent} if there
exists an integer $p \geq 1$ such that $f^p = 0$ (the zero
endomorphism).  A matrix $A \in \Mat_n(\K)$ is \emph{nilpotent} if
$A^p = 0$ for some $p \geq 1$.
\end{definition}

\begin{definition}{Index of nilpotency}{def:nilpotent-index}
\index{nilpotent!index of}\index{index of nilpotency}
If $f \in \End(E)$ is nilpotent, the \emph{index of nilpotency}
(or \emph{nilpotency index}) of~$f$ is the smallest positive
integer~$r$ such that $f^r = 0$.  We then have
\[
  f^{r-1} \neq 0 \quad\text{and}\quad f^r = 0.
\]
\end{definition}

\begin{proposition}{Basic properties of nilpotent endomorphisms}%
{prop:nilpotent-properties}
\index{nilpotent!properties}
Let $f \in \End(E)$ be nilpotent with $\dim E = n$.  Then:
\begin{enumerate}[label=(\roman*)]
  \item\label{it:nilp-eigenvalue}
    The only eigenvalue of~$f$ is $\lambda = 0$.
  \item\label{it:nilp-trace}
    $\tr(f) = 0$.
  \item\label{it:nilp-det}
    $\det(f) = 0$; in particular, $f$ is not invertible.
  \item\label{it:nilp-charpoly}
    The characteristic polynomial of~$f$ is $\chi_f(\lambda) = (-\lambda)^n$.
  \item\label{it:nilp-index-bound}
    The index of nilpotency satisfies $r \leq n$.
  \item\label{it:nilp-minpoly}
    The minimal polynomial of~$f$ is $\mu_f(\lambda) = (-\lambda)^r$,
    where $r$ is the index of nilpotency.
\end{enumerate}
\end{proposition}

\begin{proof}
\ref{it:nilp-eigenvalue}\;
Suppose $f(v) = \lambda v$ for some $v \neq 0$.  Then
$f^p(v) = \lambda^p v$, so $f^p = 0$ implies $\lambda^p v = 0$,
hence $\lambda^p = 0$ (since $v \neq 0$), so $\lambda = 0$.

\ref{it:nilp-trace}\ and \ref{it:nilp-det}\;
Since $f$ is trigonalizable (every endomorphism over~$\C$ is, and
the argument extends because all eigenvalues are~$0$), there
exists a basis in which the matrix of~$f$ is upper triangular with
all diagonal entries equal to~$0$.  Hence $\tr(f) = 0$ and
$\det(f) = 0$.

\ref{it:nilp-charpoly}\;
From the triangular form, $\chi_f(\lambda) = (0 - \lambda)^n =
(-\lambda)^n$.

\ref{it:nilp-index-bound}\;
By Cayley--Hamilton, $\chi_f(f) = 0$, i.e.\ $(-1)^n f^n = 0$, so
$f^n = 0$.  The index of nilpotency, being the smallest such
exponent, satisfies $r \leq n$.

\ref{it:nilp-minpoly}\;
Since $\chi_f(\lambda) = (-\lambda)^n$ and the minimal polynomial
divides the characteristic polynomial and has the same roots, we
have $\mu_f(\lambda) = (-\lambda)^s$ for some $1 \leq s \leq n$.
By definition, $\mu_f(f) = 0$ means $f^s = 0$, and $s$ is minimal
with this property.  Hence $s = r$.
\end{proof}

\begin{example}{Nilpotent matrices}{ex:nilpotent-matrices}
\index{nilpotent!examples}
\begin{enumerate}[label=(\alph*)]
  \item The matrix $N = \begin{pmatrix} 0 & 1 \\ 0 & 0 \end{pmatrix}$
    satisfies $N^2 = 0$, so it is nilpotent with index~$2$.

  \item The matrix
    $N = \begin{pmatrix} 0 & 1 & 0 \\ 0 & 0 & 1 \\ 0 & 0 & 0
    \end{pmatrix}$ satisfies $N^2 \neq 0$ but $N^3 = 0$, so it is
    nilpotent with index~$3$.

  \item The zero matrix $0_n$ is nilpotent with index~$1$.

  \item The matrix
    $\begin{pmatrix} 0 & 1 & 0 \\ 0 & 0 & 0 \\ 0 & 0 & 0
    \end{pmatrix}$ is nilpotent with index~$2$.
\end{enumerate}
\end{example}

\begin{proposition}{Kernel filtration of a nilpotent endomorphism}%
{prop:nilpotent-kernel-chain}
\index{nilpotent!kernel filtration}
Let $f \in \End(E)$ be nilpotent with index of nilpotency~$r$.  Then
\[
  \set{\vzero}
  \subsetneq \Ker(f)
  \subsetneq \Ker(f^2)
  \subsetneq \cdots
  \subsetneq \Ker(f^{r-1})
  \subsetneq \Ker(f^r)
  = E.
\]
In particular, each inclusion is strict, and $\dim\Ker(f^k) \geq k$
for all $1 \leq k \leq r$.
\end{proposition}

\begin{proof}
The inclusions $\Ker(f^k) \subseteq \Ker(f^{k+1})$ follow from the
fact that $f^k(v) = 0$ implies $f^{k+1}(v) = f(f^k(v)) = f(0) = 0$.
Since $f^r = 0$, we have $\Ker(f^r) = E$.  It remains to show that
each inclusion is strict.

Suppose for contradiction that $\Ker(f^k) = \Ker(f^{k+1})$ for
some $k < r$.  We claim this forces $\Ker(f^k) = \Ker(f^{k+j})$
for all $j \geq 0$.  Indeed, let $v \in \Ker(f^{k+2})$, so
$f^{k+2}(v) = 0$, meaning $f(v) \in \Ker(f^{k+1}) = \Ker(f^k)$,
hence $f^{k+1}(v) = f^k(f(v)) = 0$, so $v \in \Ker(f^{k+1}) =
\Ker(f^k)$.  By induction, $\Ker(f^k) = \Ker(f^m)$ for all
$m \geq k$.  In particular, $\Ker(f^k) = \Ker(f^r) = E$, so
$f^k = 0$, contradicting the minimality of~$r$ since $k < r$.

Since the inclusions are strict in a space of dimension~$n$, at
each step the dimension increases by at least~$1$, giving
$\dim\Ker(f^k) \geq k$.
\end{proof}


% ============================================================
\section{Invariant subspaces}\label{sec:invariant-subspaces}
% ============================================================

\begin{definition}{Invariant subspace}{def:invariant-subspace}
\index{invariant subspace}
Let $f \in \End(E)$ and let $F \subseteq E$ be a subspace.
We say that $F$ is \emph{$f$-invariant} (or \emph{invariant under
$f$}) if $f(F) \subseteq F$, i.e.\ $f(v) \in F$ for every $v \in F$.
\end{definition}

\begin{remark}{Restriction to an invariant subspace}{rem:restriction-invariant}
If $F$ is $f$-invariant, then $f$ restricts to a well-defined
endomorphism $f\restrict{F} \in \End(F)$.  This simple observation is
the key to decomposing endomorphisms into simpler pieces.
\end{remark}

\begin{proposition}{Examples of invariant subspaces}%
{prop:invariant-examples}
\index{invariant subspace!examples}
Let $f \in \End(E)$.  The following are $f$-invariant subspaces:
\begin{enumerate}[label=(\roman*)]
  \item $\set{\vzero}$ and $E$ (trivial invariant subspaces).
  \item $\Ker(f)$ and $\im(f)$.
  \item More generally, $\Ker(f^k)$ and $\im(f^k)$ for every $k \geq 1$.
  \item Every eigenspace $E_\lambda(f) = \Ker(f - \lambda\Id)$.
  \item $\Ker(P(f))$ for any polynomial $P \in \K[\lambda]$.
\end{enumerate}
\end{proposition}

\begin{proof}
We prove~(v), which implies the others as special cases.  Let
$P \in \K[\lambda]$ and $v \in \Ker(P(f))$, so $P(f)(v) = 0$.
Since $f$ commutes with $P(f)$ (both are polynomials in~$f$), we
have
\[
  P(f)\bigl(f(v)\bigr) = f\bigl(P(f)(v)\bigr) = f(0) = 0,
\]
so $f(v) \in \Ker(P(f))$.
\end{proof}

\begin{proposition}{Block-diagonal structure from invariant decomposition}%
{prop:block-diagonal}
\index{invariant subspace!block-diagonal form}
Let $E = F_1 \oplus F_2 \oplus \cdots \oplus F_s$ where each~$F_i$ is
$f$-invariant.  If $\cB_i$ is a basis of~$F_i$ and
$\cB = \cB_1 \cup \cB_2 \cup \cdots \cup \cB_s$ is the resulting
basis of~$E$, then the matrix of~$f$ in~$\cB$ is block-diagonal:
\[
  \operatorname{Mat}_\cB(f) = \begin{pmatrix}
    A_1 &        &        \\
        & A_2    &        \\
        &        & \ddots \\
        &        &        & A_s
  \end{pmatrix},
\]
where $A_i = \operatorname{Mat}_{\cB_i}(f\restrict{F_i})$.
Moreover, $\chi_f = \chi_{f\restrict{F_1}} \cdots
\chi_{f\restrict{F_s}}$.
\end{proposition}

\begin{proof}
Since each $F_i$ is $f$-invariant, $f$ maps basis vectors of
$\cB_i$ to linear combinations of vectors in $\cB_i$ only.  Hence
the matrix has the stated block-diagonal form.  The factorization
of the characteristic polynomial follows from the multiplicativity
of determinants for block-diagonal matrices.
\end{proof}


% ============================================================
\section{Generalized eigenspaces}\label{sec:generalized-eigenspaces}
% ============================================================

When the ordinary eigenspace $E_\lambda(f) = \Ker(f - \lambda\Id)$ is
too small (i.e.\ $\dim E_\lambda(f) < m_a(\lambda)$), we enlarge it
by looking at the kernels of higher powers of $f - \lambda\Id$.

\begin{definition}{Generalized eigenvector and generalized eigenspace}%
{def:generalized-eigenspace}
\index{generalized eigenvector}\index{generalized eigenspace}
\index{eigenspace!generalized}
Let $f \in \End(E)$ and $\lambda \in \K$.  A nonzero vector
$v \in E$ is a \emph{generalized eigenvector} of~$f$ for the
eigenvalue~$\lambda$ if there exists an integer $k \geq 1$ such that
\[
  (f - \lambda\Id)^k(v) = 0.
\]
The \emph{generalized eigenspace} of~$f$ for~$\lambda$ is
\[
  G_\lambda(f) \deff \bigcup_{k=1}^{\infty} \Ker\bigl(
    (f - \lambda\Id)^k
  \bigr).
\]
\end{definition}

\begin{remark}{Stabilization of the kernel chain}%
{rem:generalized-eigenspace-stabilizes}
Since $\dim E = n$ is finite, the ascending chain
$\Ker(f - \lambda\Id) \subseteq \Ker(f - \lambda\Id)^2 \subseteq
\cdots$ must stabilize at some index $k_0 \leq n$.  Therefore
\[
  G_\lambda(f) = \Ker\bigl((f - \lambda\Id)^{k_0}\bigr)
\]
for some $k_0 \leq n$, and the union in the definition is in fact
achieved at a finite stage.
\end{remark}

\begin{proposition}{Properties of generalized eigenspaces}%
{prop:generalized-eigenspace-properties}
\index{generalized eigenspace!properties}
Let $f \in \End(E)$ and suppose $\chi_f$ splits over~$\K$ as
\[
  \chi_f(\lambda) = \prod_{i=1}^{s} (\lambda_i - \lambda)^{n_i},
\]
where $\lambda_1, \ldots, \lambda_s$ are the distinct eigenvalues
and $n_i = m_a(\lambda_i)$ is the algebraic multiplicity.  Then:
\begin{enumerate}[label=(\roman*)]
  \item\label{it:gen-eig-f-inv}
    Each $G_{\lambda_i}(f)$ is $f$-invariant.
  \item\label{it:gen-eig-dim}
    $\dim G_{\lambda_i}(f) = n_i$.
  \item\label{it:gen-eig-kernel}
    $G_{\lambda_i}(f) = \Ker\bigl((f - \lambda_i\Id)^{n_i}\bigr)$.
  \item\label{it:gen-eig-restriction}
    The restriction $(f - \lambda_i\Id)\restrict{G_{\lambda_i}(f)}$
    is nilpotent.
  \item\label{it:gen-eig-direct-sum}
    The generalized eigenspaces are in direct sum:
    \[
      E = G_{\lambda_1}(f) \oplus G_{\lambda_2}(f) \oplus \cdots
        \oplus G_{\lambda_s}(f).
    \]
\end{enumerate}
\end{proposition}

\begin{proof}
\ref{it:gen-eig-f-inv}\;
This follows from 
ef{prop:prop:invariant-examples}, since
$G_\lambda(f) = \Ker\bigl((f - \lambda\Id)^k\bigr)$ for
sufficiently large~$k$, and the kernel of any polynomial in~$f$ is
$f$-invariant.

\medskip\noindent
We prove \ref{it:gen-eig-dim}, \ref{it:gen-eig-kernel}, and
\ref{it:gen-eig-direct-sum} together using the primary decomposition
theorem.

Write $\chi_f(\lambda) = \prod_{i=1}^s (\lambda_i - \lambda)^{n_i}$.
Define the polynomials $P_i(\lambda) = (\lambda_i - \lambda)^{n_i}$
and $Q_i(\lambda) = \prod_{j \neq i} (\lambda_j - \lambda)^{n_j}$.
Then $\gcd(P_i, Q_i) = 1$ since the~$\lambda_j$ are distinct.
By B\'ezout's identity in~$\K[\lambda]$, there exist polynomials
$U_i, V_i$ such that $U_i P_i + V_i Q_i = 1$.

By Cayley--Hamilton, $\chi_f(f) = 0$, i.e.\ $P_i(f)\,Q_i(f) = 0$
on $G_{\lambda_i}(f)$, and indeed
$\prod_{i=1}^s P_i(f) = 0$ on all of~$E$.

From $U_i(f) P_i(f) + V_i(f) Q_i(f) = \Id$, for any $v \in E$ we
have $v = U_i(f)P_i(f)(v) + V_i(f)Q_i(f)(v)$.  One checks that
$V_i(f)Q_i(f)(v) \in \Ker(P_i(f)) = \Ker\bigl((f -
\lambda_i\Id)^{n_i}\bigr)$ by applying $P_i(f)$ and using
$P_i(f)Q_i(f) = \chi_f(f) = 0$.  Similarly,
$U_i(f)P_i(f)(v) \in \Ker(Q_i(f))$.

Define $\pi_i = V_i(f)Q_i(f)$.  Then $\sum_{i=1}^s \pi_i = \Id$,
each $\pi_i$ maps $E$ into $\Ker(P_i(f)) = G_{\lambda_i}(f)$, and
$\pi_i \pi_j = 0$ for $i \neq j$.  This gives the direct sum
decomposition $E = \bigoplus_{i=1}^s G_{\lambda_i}(f)$.

For the dimensions, since
$\chi_{f\restrict{G_{\lambda_i}(f)}}$ divides~$\chi_f$ and
the only eigenvalue of~$f$ on $G_{\lambda_i}(f)$ is~$\lambda_i$,
we get $\chi_{f\restrict{G_{\lambda_i}(f)}}(\lambda) =
(\lambda_i - \lambda)^{d_i}$ with $d_i = \dim G_{\lambda_i}(f)$.
Since $\chi_f = \prod_i \chi_{f\restrict{G_{\lambda_i}(f)}}$
(
ef{prop:prop:block-diagonal}), comparing degrees gives $d_i = n_i$.

\ref{it:gen-eig-restriction}\;
On $G_{\lambda_i}(f)$, by definition
$(f - \lambda_i\Id)^{n_i}\restrict{G_{\lambda_i}(f)} = 0$,
so the restriction is nilpotent with index at most~$n_i$.
\end{proof}

\begin{theorem}{Primary decomposition theorem}{thm:primary-decomposition}
\index{primary decomposition theorem}\index{decomposition!primary}
Let $f \in \End(E)$ and suppose $\chi_f$ splits over~$\K$ as
$\chi_f(\lambda) = \prod_{i=1}^{s} (\lambda_i - \lambda)^{n_i}$.
Then
\[
  \boxed{E = G_{\lambda_1}(f) \oplus G_{\lambda_2}(f) \oplus \cdots
    \oplus G_{\lambda_s}(f),}
\]
where $\dim G_{\lambda_i}(f) = n_i$ and each $G_{\lambda_i}(f)$ is
$f$-invariant.  Moreover, the restriction $f\restrict{G_{\lambda_i}(f)}$
has the unique eigenvalue~$\lambda_i$.
\end{theorem}

\begin{proof}
This was established in the proof of

ef{prop:prop:generalized-eigenspace-properties}.
\end{proof}


% ============================================================
\section{Jordan blocks and Jordan matrices}\label{sec:jordan-blocks}
% ============================================================

\begin{definition}{Jordan block}{def:jordan-block}
\index{Jordan block}\index{block!Jordan}
For $\lambda \in \K$ and $k \geq 1$, the \emph{Jordan block} of
size~$k$ associated with~$\lambda$ is the $k \times k$ matrix
\[
  J_k(\lambda) \deff
  \begin{pmatrix}
    \lambda & 1       &        &         \\
            & \lambda & 1      &         \\
            &         & \ddots & \ddots  \\
            &         &        & \lambda & 1 \\
            &         &        &         & \lambda
  \end{pmatrix}
  = \lambda I_k + N_k,
\]
where $N_k$ is the $k \times k$ nilpotent matrix with $1$'s on the
superdiagonal and $0$'s elsewhere.
\end{definition}

\begin{remark}{Properties of a Jordan block}{rem:jordan-block-props}
\index{Jordan block!properties}
Let $J = J_k(\lambda)$.
\begin{enumerate}[label=(\roman*)]
  \item $\chi_J(\mu) = (\lambda - \mu)^k$.
  \item The only eigenvalue of~$J$ is~$\lambda$, with algebraic
    multiplicity~$k$ and geometric multiplicity~$1$.
  \item $J - \lambda I_k = N_k$ is nilpotent with index~$k$.
  \item $\mu_J(\mu) = (\lambda - \mu)^k$ (the minimal polynomial
    equals the characteristic polynomial).
\end{enumerate}
\end{remark}

\begin{definition}{Jordan matrix}{def:jordan-matrix}
\index{Jordan matrix}\index{matrix!Jordan}
A \emph{Jordan matrix} is a block-diagonal matrix of the form
\[
  J = \begin{pmatrix}
    J_{k_1}(\lambda_1) &                     &        \\
                        & J_{k_2}(\lambda_2) &        \\
                        &                     & \ddots \\
                        &                     &        & J_{k_m}(\lambda_m)
  \end{pmatrix},
\]
where each $J_{k_i}(\lambda_i)$ is a Jordan block.  The eigenvalues
$\lambda_1, \ldots, \lambda_m$ need not be distinct.
\end{definition}


% ============================================================
\section{The Jordan Normal Form theorem}\label{sec:jordan-theorem}
% ============================================================

\begin{theorem}{Jordan Normal Form}{thm:jordan-normal-form}
\index{Jordan normal form!theorem}
Let $E$ be a finite-dimensional vector space over an algebraically
closed field~$\K$ (e.g.\ $\K = \C$), and let $f \in \End(E)$.  Then
there exists a basis~$\cB$ of~$E$ in which the matrix of~$f$ is a
Jordan matrix:
\[
  \operatorname{Mat}_\cB(f)
  = \begin{pmatrix}
    J_{k_1}(\lambda_1) &                     &        \\
                        & J_{k_2}(\lambda_2) &        \\
                        &                     & \ddots \\
                        &                     &        & J_{k_m}(\lambda_m)
  \end{pmatrix}.
\]
Equivalently, every matrix $A \in \Mat_n(\K)$ is similar to a Jordan
matrix: $A = P J \inv{P}$ for some invertible~$P$.

Moreover, the Jordan matrix is unique up to permutation of the blocks.
\end{theorem}

The proof occupies the next section.  The overall strategy is:
\begin{enumerate}
  \item Use the primary decomposition to reduce to the case where $f$
    has a single eigenvalue.
  \item When $f$ has a single eigenvalue~$\lambda$, write
    $f = \lambda\Id + g$ where $g$ is nilpotent.
  \item Prove that every nilpotent endomorphism has a Jordan form
    (a direct sum of nilpotent Jordan blocks $J_k(0) = N_k$).
\end{enumerate}


% ============================================================
\section{Proof of the Jordan Normal Form theorem}\label{sec:jordan-proof}
% ============================================================

We begin with the nilpotent case, which is the heart of the proof.

\begin{lemma}{Jordan form for nilpotent endomorphisms}%
{lem:jordan-nilpotent}
\index{nilpotent!Jordan form}
Let $g \in \End(E)$ be nilpotent with $\dim E = n$.  Then there
exist positive integers $k_1 \geq k_2 \geq \cdots \geq k_m$ with
$k_1 + k_2 + \cdots + k_m = n$ and a basis of~$E$ in which the
matrix of~$g$ is
\[
  \begin{pmatrix}
    N_{k_1} &         &        \\
            & N_{k_2} &        \\
            &         & \ddots \\
            &         &        & N_{k_m}
  \end{pmatrix},
\]
where $N_k = J_k(0)$ is the $k \times k$ nilpotent Jordan block.
The integers $k_1, \ldots, k_m$ are uniquely determined by~$g$.
\end{lemma}

\begin{proof}
We proceed by strong induction on $n = \dim E$.

\textbf{Base case.}  If $n = 1$, then $g = 0$ and the matrix is the
$1 \times 1$ zero matrix $N_1$, which is already in Jordan form.

\textbf{Inductive step.}  Assume $n \geq 2$ and the result holds for
all spaces of dimension $< n$.  Let $r$ be the index of nilpotency
of~$g$, so $g^r = 0$ and $g^{r-1} \neq 0$.

If $r = 1$, then $g = 0$ and the matrix is the zero matrix, which is
a direct sum of $n$ copies of $N_1$.

Assume $r \geq 2$.  Consider the subspace $F = \im(g)$.  Since
$g \neq 0$, we have $F \neq \set{\vzero}$ and $\dim F < n$ (because
$g$ is not invertible).  Moreover, $F$ is $g$-invariant: for any
$v \in E$, $g(g(v)) = g^2(v) \in \im(g) = F$.

The restriction $g\restrict{F}$ is nilpotent on~$F$, and
$\dim F < n$.  By the induction hypothesis, there exists a basis of~$F$
in which $g\restrict{F}$ has the form
$\diag(N_{k_1-1}, N_{k_2-1}, \ldots, N_{k_p-1}, N_{k_{p+1}}, \ldots)$
(rearranging and shifting indices as needed).

Concretely, there exist \emph{cyclic chains} for~$g$.  Choose a vector
$v_1 \in E$ such that $g^{r-1}(v_1) \neq 0$.  Then the vectors
\[
  v_1,\; g(v_1),\; g^2(v_1),\; \ldots,\; g^{r-1}(v_1)
\]
are linearly independent (if $\sum_{j=0}^{r-1} \alpha_j g^j(v_1) = 0$,
applying $g^{r-1-j}$ for the smallest~$j$ with $\alpha_j \neq 0$ gives
a contradiction).  Write $W_1 = \Span(v_1, g(v_1), \ldots,
g^{r-1}(v_1))$.  In the basis
$(g^{r-1}(v_1), g^{r-2}(v_1), \ldots, g(v_1), v_1)$, the matrix
of $g\restrict{W_1}$ is $N_r$.

We now construct a complement.  Set
$d_j = \dim\Ker(g^j)$ for $0 \leq j \leq r$.  Define
$m_j = 2d_j - d_{j-1} - d_{j+1}$ for $1 \leq j \leq r-1$ and
$m_r = 2d_r - d_{r-1} - d_r = d_r - d_{r-1}$.  We claim that~$m_j$
equals the number of Jordan blocks of size~$j$.

More precisely, we construct the full Jordan basis as follows.
For each $j$ from $r$ down to~$1$, choose vectors $v$ such that
$g^{j-1}(v) \neq 0$ but $g^{j}(v) = 0$, and such that
$v, g(v), \ldots, g^{j-1}(v)$ are linearly independent of the vectors
already chosen.  The number of such chains of length~$j$ is
$m_j = d_j - d_{j-1} - (\text{number of chain elements of length}
> j\text{ lying in }\Ker(g^j)\setminus\Ker(g^{j-1}))$.

To be fully explicit, define $\ell_j = d_j - d_{j-1}$ (with $d_0 = 0$).
One can show that $\ell_1 \geq \ell_2 \geq \cdots \geq \ell_r \geq 1$
and that $\ell_j$ equals the number of Jordan blocks of size $\geq j$.
Hence the number of blocks of size exactly~$j$ is
$\ell_j - \ell_{j+1}$ (where $\ell_{r+1} = 0$).

We verify this by induction.  The subspace $W_1$ Spanned by the first
chain has the property that $g\restrict{W_1}$ corresponds to $N_r$,
and $E = W_1 \oplus W_1'$ where $W_1'$ is a $g$-stable complement.
Such a complement exists: choose a complement $C$ of $\Ker(g^{r-1})$
inside $\Ker(g^r) = E$, so $\dim C = \ell_r$.  For each basis
vector~$c$ of~$C$, the chain $c, g(c), \ldots, g^{r-1}(c)$ Spans
a cyclic subspace isomorphic to $N_r$.  We can choose these chains so
that their union is linearly independent.  After removing these chains,
the remaining space (a suitable complement) has a nilpotent restriction
of strictly smaller dimension, and the induction hypothesis applies.

This procedure produces a basis of~$E$ consisting of cyclic chains,
and in this basis the matrix of~$g$ is block-diagonal with nilpotent
Jordan blocks.

\textbf{Uniqueness.}  The sizes of the Jordan blocks are determined
by the sequence $d_j = \dim\Ker(g^j)$: the number of blocks of
size~$\geq j$ is $d_j - d_{j-1}$, and the number of blocks of size
exactly~$j$ is $(d_j - d_{j-1}) - (d_{j+1} - d_j)$.  Since the
dimensions $d_j$ are invariants of~$g$, the block sizes are uniquely
determined.
\end{proof}

\begin{proof}[Proof of 
ef{thm:thm:jordan-normal-form}]
\index{Jordan normal form!proof}

\textbf{Existence.}  By the primary decomposition theorem
(
ef{thm:thm:primary-decomposition}),
\[
  E = G_{\lambda_1}(f) \oplus G_{\lambda_2}(f) \oplus \cdots
    \oplus G_{\lambda_s}(f),
\]
where $\lambda_1, \ldots, \lambda_s$ are the distinct eigenvalues
of~$f$ and $\dim G_{\lambda_i}(f) = n_i = m_a(\lambda_i)$.

On each $G_{\lambda_i}(f)$, define
$g_i = (f - \lambda_i\Id)\restrict{G_{\lambda_i}(f)}$, which is
nilpotent.  By 
ef{lem:lem:jordan-nilpotent}, there is a basis of
$G_{\lambda_i}(f)$ in which $g_i$ has the form
$\diag(N_{k_1^{(i)}}, \ldots, N_{k_{m_i}^{(i)}})$.

Since $f\restrict{G_{\lambda_i}(f)} = \lambda_i\Id + g_i$, the
matrix of $f\restrict{G_{\lambda_i}(f)}$ in this basis is
\[
  \diag\bigl(J_{k_1^{(i)}}(\lambda_i), \ldots,
    J_{k_{m_i}^{(i)}}(\lambda_i)\bigr).
\]
Concatenating the bases of all generalized eigenspaces gives a basis
of~$E$ in which $f$ has a Jordan matrix.

\textbf{Uniqueness.}  On each generalized eigenspace
$G_{\lambda_i}(f)$, the Jordan block sizes are determined by
the dimensions $\dim\Ker\bigl((f - \lambda_i\Id)^k\bigr)$
by 
ef{lem:lem:jordan-nilpotent}.  Since the eigenvalues, their
algebraic multiplicities, and the kernel dimensions are similarity
invariants, the Jordan form is unique up to the ordering of blocks.
\end{proof}


% ============================================================
\section{Computing the Jordan form}\label{sec:computing-jordan}
% ============================================================

Given a matrix $A \in \Mat_n(\K)$, the following algorithm computes
its Jordan normal form.

\begin{enumerate}[label=\textbf{Step \arabic*:},leftmargin=3cm]
  \item \textbf{Find the eigenvalues.}
    Compute $\chi_A(\lambda) = \det(A - \lambda I)$ and factor it:
    $\chi_A(\lambda) = \prod_{i=1}^s (\lambda_i - \lambda)^{n_i}$.

  \item \textbf{For each eigenvalue~$\lambda_i$, compute the kernel
    dimensions.}
    Compute $d_k^{(i)} = \dim\Ker\bigl((A - \lambda_i I)^k\bigr)$
    for $k = 1, 2, 3, \ldots$ until $d_k^{(i)} = n_i$ (which must
    happen for some $k \leq n_i$).

  \item \textbf{Determine the Jordan block sizes.}
    For each eigenvalue~$\lambda_i$:
    \begin{itemize}
      \item The number of Jordan blocks for $\lambda_i$ is
        $d_1^{(i)} = \dim\Ker(A - \lambda_i I)$ (the geometric
        multiplicity).
      \item The number of blocks of size $\geq k$ is
        $d_k^{(i)} - d_{k-1}^{(i)}$.
      \item The number of blocks of size exactly~$k$ is
        $(d_k^{(i)} - d_{k-1}^{(i)}) - (d_{k+1}^{(i)} - d_k^{(i)})
        = 2d_k^{(i)} - d_{k-1}^{(i)} - d_{k+1}^{(i)}$.
    \end{itemize}

  \item \textbf{Assemble the Jordan matrix.}
    Write $J = \diag\bigl(J_{k_1}(\lambda_1), J_{k_2}(\lambda_2),
    \ldots\bigr)$.
\end{enumerate}

\begin{remark}{Finding the transition matrix}{rem:jordan-transition}
\index{Jordan normal form!transition matrix}
To find the invertible matrix~$P$ such that $A = P J \inv{P}$, one
must compute a \emph{Jordan basis}: for each Jordan block of
size~$k$ associated with eigenvalue~$\lambda$, find a chain of
generalized eigenvectors $v_1, v_2, \ldots, v_k$ satisfying
\[
  (A - \lambda I)v_1 = 0, \quad
  (A - \lambda I)v_2 = v_1, \quad
  \ldots, \quad
  (A - \lambda I)v_k = v_{k-1}.
\]
The columns of~$P$ are then these generalized eigenvectors
$(v_k, v_{k-1}, \ldots, v_1)$ (note the reversed order: $v_k$ comes
first so that $AP = PJ$).
\end{remark}


% ============================================================
\section{Worked examples}\label{sec:jordan-examples}
% ============================================================

\begin{example}{Jordan form of a $3 \times 3$ matrix}{ex:jordan-3x3}
\index{Jordan normal form!example, $3\times 3$}
Let
\[
  A = \begin{pmatrix} 2 & 1 & 0 \\ 0 & 2 & 1 \\ 0 & 0 & 2
  \end{pmatrix}.
\]

\textbf{Step 1.}  $\chi_A(\lambda) = (2 - \lambda)^3$, so $\lambda = 2$
is the only eigenvalue with $n_1 = 3$.

\textbf{Step 2.}  Compute the kernel dimensions.
\[
  A - 2I = \begin{pmatrix} 0 & 1 & 0 \\ 0 & 0 & 1 \\ 0 & 0 & 0
  \end{pmatrix}, \quad
  (A - 2I)^2 = \begin{pmatrix} 0 & 0 & 1 \\ 0 & 0 & 0 \\ 0 & 0 & 0
  \end{pmatrix}, \quad
  (A - 2I)^3 = 0.
\]
So $d_1 = \dim\Ker(A - 2I) = 1$, $d_2 = \dim\Ker(A - 2I)^2 = 2$,
$d_3 = 3$.

\textbf{Step 3.}  Number of blocks of size~$\geq 1$: $d_1 - d_0 = 1$.
Number of blocks of size~$\geq 2$: $d_2 - d_1 = 1$.
Number of blocks of size~$\geq 3$: $d_3 - d_2 = 1$.
So there is exactly one block of size~$3$.

\textbf{Step 4.}  The Jordan form is
\[
  J = J_3(2) = \begin{pmatrix} 2 & 1 & 0 \\ 0 & 2 & 1 \\ 0 & 0 & 2
  \end{pmatrix}.
\]
Indeed, $A$ is already in Jordan form.
\end{example}

\begin{example}{Jordan form of a $4 \times 4$ matrix}{ex:jordan-4x4}
\index{Jordan normal form!example, $4\times 4$}
Let
\[
  A = \begin{pmatrix}
    3 & 1 & 0 & 0 \\
    0 & 3 & 0 & 0 \\
    0 & 0 & 3 & 1 \\
    0 & 0 & 0 & 5
  \end{pmatrix}.
\]

\textbf{Step 1.}  $\chi_A(\lambda) = (3 - \lambda)^3(5 - \lambda)$.
Eigenvalues: $\lambda_1 = 3$ with $n_1 = 3$, and $\lambda_2 = 5$
with $n_2 = 1$.

\textbf{Step 2.}  For $\lambda_1 = 3$:
\[
  A - 3I = \begin{pmatrix}
    0 & 1 & 0 & 0 \\
    0 & 0 & 0 & 0 \\
    0 & 0 & 0 & 1 \\
    0 & 0 & 0 & 2
  \end{pmatrix}, \quad
  (A - 3I)^2 = \begin{pmatrix}
    0 & 0 & 0 & 0 \\
    0 & 0 & 0 & 0 \\
    0 & 0 & 0 & 2 \\
    0 & 0 & 0 & 4
  \end{pmatrix}.
\]
$d_1 = \dim\Ker(A - 3I) = 2$ (the kernel is Spanned by
$e_1$ and $e_3$), $d_2 = \dim\Ker(A-3I)^2 = 2$.

Wait---we need $d_k = \dim\Ker(A-3I)^k$ restricted to the generalized
eigenspace, or equivalently the rank of $(A - 3I)^k$.
Since $d_2 = 2$ and we need $d_k = 3 = n_1$, let us recompute.

Actually, $\rank(A - 3I) = 2$ (two independent rows), so
$d_1 = 4 - 2 = 2$.  For $(A - 3I)^2$: $\rank = 1$, so
$d_2 = 4 - 1 = 3$.

Number of blocks for $\lambda_1 = 3$: $d_1 = 2$ blocks total.
Blocks of size $\geq 2$: $d_2 - d_1 = 1$.
Blocks of size $\geq 3$: $d_3 - d_2 = 0$ (since $d_3 = 3 = d_2$
for the generalized eigenspace dimension).
So: one block of size~$2$ (from $1$ of size $\geq 2$, minus $0$ of
size $\geq 3$), and one block of size~$1$ (from $2$ total, minus $1$
of size~$\geq 2$).

For $\lambda_2 = 5$: one block of size~$1$ (the algebraic multiplicity
is~$1$).

\textbf{Step 3.}  The Jordan form is
\[
  J = \begin{pmatrix}
    3 & 1 & 0 & 0 \\
    0 & 3 & 0 & 0 \\
    0 & 0 & 3 & 0 \\
    0 & 0 & 0 & 5
  \end{pmatrix}
  = \diag\bigl(J_2(3),\; J_1(3),\; J_1(5)\bigr).
\]

\textbf{Jordan basis.}  For the block $J_2(3)$: we need $v_1 \in
\Ker(A - 3I)$ and $v_2$ with $(A - 3I)v_2 = v_1$.  Take
$v_1 = e_1$, then $(A - 3I)v_2 = e_1$ gives $v_2 = e_2$ (since
$(A-3I)e_2 = e_1$).  For $J_1(3)$: take $v_3 = e_3 \in \Ker(A-3I)$.
For $J_1(5)$: solve $(A - 5I)v_4 = 0$; we find
$v_4 = \begin{psmallmatrix} -1\\0\\-1\\2 \end{psmallmatrix}$
(after row reduction).

The transition matrix is $P = (v_2 \mid v_1 \mid v_3 \mid v_4)$:
\[
  P = \begin{pmatrix}
    0 & 1 & 0 & -1 \\
    1 & 0 & 0 & 0 \\
    0 & 0 & 1 & -1 \\
    0 & 0 & 0 & 2
  \end{pmatrix}, \qquad
  A = P\,J\,\inv{P}.
\]
One may verify that $AP = PJ$.
\end{example}

\begin{example}{A non-trivial $4 \times 4$ Jordan form}{ex:jordan-4x4-nontrivial}
Consider
\[
  A = \begin{pmatrix}
    5 & 4 & 2 & 1 \\
    0 & 1 & -1 & -1 \\
    -1 & -1 & 3 & 0 \\
    1 & 1 & -1 & 2
  \end{pmatrix}.
\]
One computes $\chi_A(\lambda) = (\lambda - 1)^2(\lambda - 3)^2$.

For $\lambda_1 = 1$: $\rank(A - I) = 2$, so $d_1 = 4 - 2 = 2 = n_1$.
Since $d_1 = n_1$, the generalized eigenspace equals the eigenspace.
Thus there are two blocks of size~$1$: $J_1(1) \oplus J_1(1)$.

For $\lambda_2 = 3$: $\rank(A - 3I) = 3$, so
$d_1 = 4 - 3 = 1$.  Since $d_1 = 1 < 2 = n_2$, we need to go further:
$\rank(A - 3I)^2 = 2$, so $d_2 = 4 - 2 = 2 = n_2$.
Number of blocks: $d_1 = 1$.  Size $\geq 2$: $d_2 - d_1 = 1$.
So there is one block of size~$2$: $J_2(3)$.

The Jordan form is:
\[
  J = \diag\bigl(J_1(1),\; J_1(1),\; J_2(3)\bigr)
  = \begin{pmatrix}
    1 & 0 & 0 & 0 \\
    0 & 1 & 0 & 0 \\
    0 & 0 & 3 & 1 \\
    0 & 0 & 0 & 3
  \end{pmatrix}.
\]
\end{example}


% ============================================================
\section{Jordan form and the minimal polynomial}%
\label{sec:jordan-minpoly}
% ============================================================

The Jordan form makes the relationship between a matrix and its minimal
polynomial completely transparent.

\begin{theorem}{Minimal polynomial from the Jordan form}%
{thm:jordan-minpoly}
\index{minimal polynomial!and Jordan form}
Let $A \in \Mat_n(\K)$ with Jordan form
$J = \diag\bigl(J_{k_1}(\lambda_1), \ldots, J_{k_m}(\lambda_m)\bigr)$,
where $\lambda_1, \ldots, \lambda_s$ are the distinct eigenvalues.
For each distinct eigenvalue~$\lambda_i$, let $r_i$ denote the
\emph{largest} Jordan block size associated with~$\lambda_i$.  Then
the minimal polynomial of~$A$ is
\[
  \boxed{\mu_A(\lambda) = \prod_{i=1}^{s}
    (\lambda - \lambda_i)^{r_i}.}
\]
\end{theorem}

\begin{proof}
A polynomial $P(\lambda)$ satisfies $P(A) = 0$ if and only if
$P(J) = 0$ (since $A$ and $J$ are similar).  Moreover,
$P(J) = \diag\bigl(P(J_{k_1}(\lambda_1)), \ldots,
P(J_{k_m}(\lambda_m))\bigr)$, so $P(J) = 0$ if and only if
$P(J_{k_j}(\lambda_j)) = 0$ for every block.

For a single Jordan block $J_k(\lambda)$, one computes
\[
  P(J_k(\lambda)) = \begin{pmatrix}
    P(\lambda) & P'(\lambda) & \frac{P''(\lambda)}{2!} & \cdots
      & \frac{P^{(k-1)}(\lambda)}{(k-1)!} \\
    & P(\lambda) & P'(\lambda) & \cdots
      & \frac{P^{(k-2)}(\lambda)}{(k-2)!} \\
    & & \ddots & \ddots & \vdots \\
    & & & P(\lambda) & P'(\lambda) \\
    & & & & P(\lambda)
  \end{pmatrix}.
\]
Hence $P(J_k(\lambda)) = 0$ if and only if
$P(\lambda) = P'(\lambda) = \cdots = P^{(k-1)}(\lambda) = 0$,
which holds if and only if $(\lambda_0 - \lambda)^k \mid P(\lambda)$
where $\lambda_0 = \lambda$ is the eigenvalue.

Therefore, $P(J) = 0$ if and only if $(\lambda_i - x)^{r_i}$ divides
$P(x)$ for every distinct eigenvalue~$\lambda_i$, where $r_i$ is the
largest block for~$\lambda_i$.  The monic polynomial of smallest degree
with this property is $\mu_A(\lambda) = \prod_{i=1}^s
(\lambda - \lambda_i)^{r_i}$.
\end{proof}

\begin{corollary}{Diagonalizability via the Jordan form}%
{cor:diag-jordan}
\index{diagonalizable!and Jordan form}
A matrix is diagonalizable if and only if its Jordan form is diagonal,
which happens if and only if all Jordan blocks have size~$1$.
Equivalently, the minimal polynomial has no repeated roots.
\end{corollary}

\begin{proof}
All blocks have size~$1$ if and only if $r_i = 1$ for each
eigenvalue, if and only if $\mu_A$ has only simple roots, if and only
if $A$ is diagonalizable.
\end{proof}


% ============================================================
\section{Cayley--Hamilton revisited}\label{sec:cayley-hamilton-jordan}
% ============================================================

The Jordan form gives an elegant proof of the Cayley--Hamilton theorem.

\begin{theorem}{Cayley--Hamilton theorem}{thm:cayley-hamilton-jordan}
\index{Cayley--Hamilton theorem!proof via Jordan form}
Let $A \in \Mat_n(\K)$ with $\K$ algebraically closed.  Then
$\chi_A(A) = 0$.
\end{theorem}

\begin{proof}
Let $J = P^{-1}AP$ be the Jordan form.  Then
$\chi_A(A) = P\,\chi_A(J)\,P^{-1} = P\,\chi_J(J)\,P^{-1}$
(since $\chi_A = \chi_J$).

Write $J = \diag(J_{k_1}(\lambda_1), \ldots, J_{k_m}(\lambda_m))$.
Then $\chi_J(\lambda) = \prod_{j=1}^m (\lambda_j - \lambda)^{k_j}$,
so
\[
  \chi_J(J) = \diag\bigl(\chi_J(J_{k_1}(\lambda_1)), \ldots,
    \chi_J(J_{k_m}(\lambda_m))\bigr).
\]
For each block $J_{k_j}(\lambda_j)$, the factor
$(\lambda_j - \lambda)^{k_j}$ divides $\chi_J(\lambda)$, hence
$\chi_J(J_{k_j}(\lambda_j)) = 0$ by the analysis in

ef{thm:thm:jordan-minpoly}.

Therefore $\chi_J(J) = 0$, and so $\chi_A(A) = 0$.
\end{proof}


% ============================================================
\section{Matrix exponential}\label{sec:matrix-exponential}
% ============================================================

One of the most important applications of the Jordan form is the
computation of the \emph{matrix exponential}, which plays a central
role in the theory of linear differential equations.

\begin{definition}{Matrix exponential}{def:matrix-exponential}
\index{matrix exponential}\index{exponential!matrix}
For $A \in \Mat_n(\C)$, the \emph{matrix exponential} of~$A$ is
defined by the power series
\[
  e^A \deff \exp(A) \deff \sum_{k=0}^{\infty} \frac{A^k}{k!}
  = I + A + \frac{A^2}{2!} + \frac{A^3}{3!} + \cdots.
\]
This series converges absolutely for every $A \in \Mat_n(\C)$.
\end{definition}

\begin{proposition}{Properties of the matrix exponential}%
{prop:matrix-exp-properties}
\index{matrix exponential!properties}
\begin{enumerate}[label=(\roman*)]
  \item $e^{0} = I_n$.
  \item If $AB = BA$, then $e^{A+B} = e^A e^B$.
  \item $e^A$ is always invertible, with $\inv{(e^A)} = e^{-A}$.
  \item $e^{PAP^{-1}} = P\,e^A\,P^{-1}$ for any invertible~$P$.
  \item $\det(e^A) = e^{\tr(A)}$.
\end{enumerate}
\end{proposition}

\begin{proof}
(i) is clear.  (ii) follows by expanding the product of power series,
using $AB = BA$ to collect terms as in the scalar case.  (iii) follows
from (ii) with $B = -A$.  (iv) follows from
$(PAP^{-1})^k = PA^kP^{-1}$.

For~(v), if $A = PJP^{-1}$ is in Jordan form, then $e^A = Pe^JP^{-1}$
and $\det(e^A) = \det(e^J)$.  For a Jordan block $J_k(\lambda)$,
the matrix $e^{J_k(\lambda)}$ is upper triangular with $e^\lambda$ on
the diagonal, so $\det(e^{J_k(\lambda)}) = e^{k\lambda}$.
Thus $\det(e^J) = \prod_j e^{k_j\lambda_j} = e^{\sum_j k_j\lambda_j}
= e^{\tr(J)} = e^{\tr(A)}$.
\end{proof}

\begin{proposition}{Exponential of a Jordan block}%
{prop:exp-jordan-block}
\index{matrix exponential!of Jordan block}
For a Jordan block $J_k(\lambda) = \lambda I_k + N_k$, since
$\lambda I_k$ and $N_k$ commute:
\[
  e^{J_k(\lambda)} = e^{\lambda I_k}\,e^{N_k}
  = e^{\lambda}\,e^{N_k}.
\]
Moreover, since $N_k^k = 0$:
\[
  e^{N_k} = I_k + N_k + \frac{N_k^2}{2!} + \cdots
    + \frac{N_k^{k-1}}{(k-1)!}.
\]
Explicitly:
\[
  e^{J_k(\lambda)} = e^{\lambda}
  \begin{pmatrix}
    1      & 1      & \frac{1}{2!}  & \cdots & \frac{1}{(k-1)!} \\
    0      & 1      & 1             & \cdots & \frac{1}{(k-2)!} \\
    \vdots &        & \ddots        & \ddots & \vdots            \\
    0      & \cdots &               & 1      & 1                 \\
    0      & \cdots &               & 0      & 1
  \end{pmatrix}.
\]
More generally, for the \emph{matrix exponential of~$tA$}
(with $t \in \R$):
\[
  e^{tJ_k(\lambda)} = e^{\lambda t}
  \begin{pmatrix}
    1      & t      & \frac{t^2}{2!}  & \cdots & \frac{t^{k-1}}{(k-1)!} \\
    0      & 1      & t               & \cdots & \frac{t^{k-2}}{(k-2)!} \\
    \vdots &        & \ddots          & \ddots & \vdots                  \\
    0      & \cdots &                 & 1      & t                       \\
    0      & \cdots &                 & 0      & 1
  \end{pmatrix}.
\]
\end{proposition}

\begin{proof}
Since $\lambda I_k$ and $N_k$ commute,
$e^{J_k(\lambda)} = e^{\lambda I_k + N_k} = e^{\lambda I_k} e^{N_k}
= e^\lambda I_k \cdot e^{N_k}$.  The entries of $e^{N_k}$ follow
from $(N_k^j)_{pq} = \delta_{p+j, q}$ (the matrix with $1$'s on the
$j$-th superdiagonal), giving
$(e^{N_k})_{pq} = \frac{1}{(q-p)!}$ for $q \geq p$ and $0$
otherwise.  The formula for $e^{tJ_k(\lambda)}$ is analogous, replacing
$N_k$ by $tN_k$.
\end{proof}

\begin{example}{Computing $e^{tA}$}{ex:matrix-exp}
\index{matrix exponential!example}
Let $A = \begin{pmatrix} 2 & 1 \\ 0 & 2 \end{pmatrix} = J_2(2)$.
Then
\[
  e^{tA} = e^{2t}
  \begin{pmatrix} 1 & t \\ 0 & 1 \end{pmatrix}
  = \begin{pmatrix} e^{2t} & te^{2t} \\ 0 & e^{2t} \end{pmatrix}.
\]
\end{example}


% ============================================================
\section{Applications}\label{sec:jordan-applications}
% ============================================================

% ------------------------------------------------------------
\subsection{Systems of linear ODEs}\label{subsec:jordan-ode}
% ------------------------------------------------------------

\index{differential equations!and Jordan form}

Consider the linear system of ordinary differential equations
\begin{equation}\label{eq:ode-system}
  \vec{x}'(t) = A\,\vec{x}(t), \qquad \vec{x}(0) = \vec{x}_0,
\end{equation}
where $A \in \Mat_n(\R)$ (or $\Mat_n(\C)$) and
$\vec{x}(t) \in \R^n$ (or $\C^n$).

The solution is
$\vec{x}(t) = e^{tA}\vec{x}_0$.
When $A$ is diagonalizable, $A = P\diag(\lambda_1,\ldots,\lambda_n)
\inv{P}$ and the solution is a superposition of pure exponentials
$e^{\lambda_i t}$.  When $A$ is \emph{not} diagonalizable, the Jordan
form tells us exactly what happens.

\begin{proposition}{Solution structure via Jordan form}%
{prop:ode-jordan}
\index{differential equations!non-diagonalizable case}
If $A = PJ\inv{P}$ with Jordan form~$J$, then
$e^{tA} = P\,e^{tJ}\,\inv{P}$ and
$\vec{x}(t) = P\,e^{tJ}\,\inv{P}\vec{x}_0$.

For a Jordan block $J_k(\lambda)$, the corresponding components of
the solution involve terms of the form
\[
  t^j e^{\lambda t}, \qquad 0 \leq j \leq k-1.
\]
In particular, non-trivial Jordan blocks (size $> 1$) produce
\emph{polynomial-exponential} solutions (polynomial times
exponential), not pure exponentials.
\end{proposition}

\begin{example}{A $3 \times 3$ system of ODEs}{ex:ode-3x3}
Consider $\vec{x}'(t) = A\vec{x}(t)$ with
\[
  A = \begin{pmatrix} 1 & 1 & 0 \\ 0 & 1 & 1 \\ 0 & 0 & 1
  \end{pmatrix} = J_3(1).
\]
Then
\[
  e^{tA} = e^{t}
  \begin{pmatrix}
    1 & t & \tfrac{t^2}{2} \\
    0 & 1 & t \\
    0 & 0 & 1
  \end{pmatrix}.
\]
With initial condition
$\vec{x}_0 = \begin{psmallmatrix} 1\\0\\1 \end{psmallmatrix}$:
\[
  \vec{x}(t) = e^{t}
  \begin{pmatrix}
    1 + \tfrac{t^2}{2} \\
    t \\
    1
  \end{pmatrix}.
\]
Note the appearance of the polynomial factor $1 + t^2/2$ in the first
component, which is characteristic of a $3\times 3$ Jordan block.
\end{example}

% ------------------------------------------------------------
\subsection{Matrix functions via Jordan form}%
\label{subsec:matrix-functions}
% ------------------------------------------------------------

\index{matrix function}

The Jordan form allows us to define and compute general matrix
functions.  If $h\colon \C \to \C$ is a function that is sufficiently
smooth (or analytic) in a neighbourhood of each eigenvalue, we can
define $h(A)$ as follows.

\begin{definition}{Matrix function via Jordan form}%
{def:matrix-function}
Let $A \in \Mat_n(\C)$ with $A = PJ\inv{P}$ and
$J = \diag(J_{k_1}(\lambda_1), \ldots, J_{k_m}(\lambda_m))$.
Define
\[
  h(A) \deff P\,\diag\bigl(h(J_{k_1}(\lambda_1)), \ldots,
    h(J_{k_m}(\lambda_m))\bigr)\,\inv{P},
\]
where
\[
  h(J_k(\lambda)) \deff
  \begin{pmatrix}
    h(\lambda) & h'(\lambda) & \frac{h''(\lambda)}{2!} & \cdots
      & \frac{h^{(k-1)}(\lambda)}{(k-1)!} \\
    & h(\lambda) & h'(\lambda) & \cdots
      & \frac{h^{(k-2)}(\lambda)}{(k-2)!} \\
    & & \ddots & \ddots & \vdots \\
    & & & h(\lambda) & h'(\lambda) \\
    & & & & h(\lambda)
  \end{pmatrix}.
\]
\end{definition}

\begin{example}{Computing $A^k$ and $\sqrt{A}$}{ex:matrix-functions}
For $A = \begin{psmallmatrix} 4 & 1 \\ 0 & 4 \end{psmallmatrix}
= J_2(4)$ and $h(\lambda) = \lambda^k$:
\[
  A^k = \begin{pmatrix} 4^k & k\cdot 4^{k-1} \\ 0 & 4^k
  \end{pmatrix}.
\]
For $h(\lambda) = \sqrt{\lambda}$ (choosing the principal square root):
\[
  \sqrt{A} = \begin{pmatrix} 2 & \tfrac{1}{4} \\ 0 & 2
  \end{pmatrix},
\]
since $h(4) = 2$ and $h'(4) = \frac{1}{2\sqrt{4}} = \frac{1}{4}$.
\end{example}


% ============================================================
\section{Exercises}\label{sec:jordan-exercises}
% ============================================================

\begin{exercise}{Nilpotent matrices \basic}{exo:nilpotent-basic}
Determine whether each matrix is nilpotent.  If so, find the index of
nilpotency.
\begin{enumerate}[label=(\alph*)]
  \item $A = \begin{pmatrix} 0 & 2 & 1 \\ 0 & 0 & -3 \\ 0 & 0 & 0
    \end{pmatrix}$
  \item $B = \begin{pmatrix} 1 & 1 \\ 0 & 1 \end{pmatrix}$
  \item $C = \begin{pmatrix} 0 & 0 & 0 \\ 1 & 0 & 0 \\ 0 & 1 & 0
    \end{pmatrix}$
\end{enumerate}
\end{exercise}

\begin{exercise}{Properties of nilpotent endomorphisms \basic}%
{exo:nilpotent-properties}
Let $f \in \End(E)$ be nilpotent.  Show that:
\begin{enumerate}[label=(\alph*)]
  \item $\Id + f$ is invertible.
    \emph{Hint:} consider $(\Id + f)(\Id - f + f^2 - \cdots)$.
  \item If $g \in \End(E)$ is nilpotent and $fg = gf$, then $f + g$
    is nilpotent.
\end{enumerate}
\end{exercise}

\begin{exercise}{Jordan form of $2 \times 2$ matrices \basic}%
{exo:jordan-2x2}
Find the Jordan normal form of each matrix over~$\C$:
\begin{enumerate}[label=(\alph*)]
  \item $A = \begin{pmatrix} 3 & 1 \\ 0 & 3 \end{pmatrix}$
  \item $B = \begin{pmatrix} 0 & -1 \\ 1 & 0 \end{pmatrix}$
  \item $C = \begin{pmatrix} 5 & 3 \\ -3 & -1 \end{pmatrix}$
\end{enumerate}
\end{exercise}

\begin{exercise}{Jordan form of a $3 \times 3$ matrix \intermediate}%
{exo:jordan-3x3}
Find the Jordan normal form and a Jordan basis for
\[
  A = \begin{pmatrix} 2 & 1 & 0 \\ 0 & 2 & 0 \\ 0 & 0 & 3
  \end{pmatrix}.
\]
\end{exercise}

\begin{exercise}{A $4 \times 4$ matrix with repeated eigenvalues
\intermediate}{exo:jordan-4x4-repeated}
\index{Jordan normal form!exercise}
Find the Jordan normal form of
\[
  A = \begin{pmatrix}
    2 & 1 & 0 & 0 \\
    0 & 2 & 1 & 0 \\
    0 & 0 & 2 & 0 \\
    0 & 0 & 0 & 2
  \end{pmatrix}.
\]
How many distinct Jordan forms are possible for a $4 \times 4$ matrix
whose only eigenvalue is~$2$?  List them all.
\end{exercise}

\begin{exercise}{Uniqueness via kernel dimensions \intermediate}%
{exo:jordan-kernel-dimensions}
Let $A, B \in \Mat_n(\C)$ be two matrices with the same characteristic
polynomial.  Show that $A$ and $B$ are similar if and only if
$\rank(A - \lambda I)^k = \rank(B - \lambda I)^k$ for every
eigenvalue~$\lambda$ and every $k \geq 1$.
\end{exercise}

\begin{exercise}{Minimal polynomial from Jordan form \intermediate}%
{exo:minpoly-jordan}
For each of the following Jordan matrices, determine the minimal
polynomial:
\begin{enumerate}[label=(\alph*)]
  \item $J = \diag(J_3(2), J_1(2))$
  \item $J = \diag(J_2(1), J_2(3))$
  \item $J = \diag(J_2(0), J_2(0), J_1(0))$
\end{enumerate}
\end{exercise}

\begin{exercise}{Matrix exponential \intermediate}{exo:matrix-exp}
\index{matrix exponential!exercise}
Compute $e^{tA}$ for:
\begin{enumerate}[label=(\alph*)]
  \item $A = \begin{pmatrix} 0 & 1 \\ 0 & 0 \end{pmatrix}$
  \item $A = \begin{pmatrix} 1 & 1 & 0 \\ 0 & 1 & 1 \\ 0 & 0 & 1
    \end{pmatrix}$
  \item $A = \begin{pmatrix} 2 & 1 & 0 & 0 \\ 0 & 2 & 0 & 0 \\
    0 & 0 & 3 & 1 \\ 0 & 0 & 0 & 3 \end{pmatrix}$
\end{enumerate}
\end{exercise}

\begin{exercise}{Solving a non-diagonalizable ODE system
\intermediate}{exo:ode-jordan}
\index{differential equations!exercise}
Solve the initial value problem
\[
  \vec{x}'(t) = \begin{pmatrix} 3 & 1 \\ 0 & 3 \end{pmatrix}
  \vec{x}(t), \qquad
  \vec{x}(0) = \begin{pmatrix} 1 \\ 2 \end{pmatrix}.
\]
\end{exercise}

\begin{exercise}{Cayley--Hamilton via Jordan form \challenging}%
{exo:cayley-hamilton-jordan}
Let $A \in \Mat_n(\C)$ with Jordan form $J$.
\begin{enumerate}[label=(\alph*)]
  \item Compute $\chi_A(J_{k}(\lambda))$ directly, showing each
    entry is zero.
  \item Deduce the Cayley--Hamilton theorem for matrices over~$\C$.
  \item Explain why the result extends to matrices over any field~$\K$.
    \emph{Hint:} use the fact that $\chi_A(A) = 0$ is a polynomial
    identity in the entries of~$A$.
\end{enumerate}
\end{exercise}

\begin{exercise}{Powers of nilpotent-plus-scalar matrices
\challenging}{exo:powers-jordan}
Let $A = \lambda I + N$ where $N$ is nilpotent with index~$r$.
\begin{enumerate}[label=(\alph*)]
  \item Show that for all $m \geq 0$,
    \[
      A^m = \sum_{j=0}^{r-1} \binom{m}{j} \lambda^{m-j} N^j.
    \]
  \item Deduce the entries of $J_k(\lambda)^m$ explicitly.
  \item Use part~(a) to give another proof that if $A$ is nilpotent
    and $\tr(A^k) = 0$ for all $k \geq 1$, then $A$ is nilpotent.
\end{enumerate}
\end{exercise}

\begin{exercise}{Jordan form and commuting matrices \challenging}%
{exo:commuting-jordan}
\index{commuting matrices}
Let $A \in \Mat_n(\C)$.
\begin{enumerate}[label=(\alph*)]
  \item Show that $B$ commutes with $J_k(\lambda)$ if and only if $B$
    is an upper triangular Toeplitz matrix (i.e.\ $B_{ij}$ depends only
    on $j - i$).
  \item Deduce that if $A$ has $n$ distinct eigenvalues, then every
    matrix commuting with~$A$ is a polynomial in~$A$.
  \item Give an example showing this fails when $A$ has repeated
    eigenvalues with multiple Jordan blocks.
\end{enumerate}
\end{exercise}

\begin{exercise}{Logarithm of a matrix \challenging}{exo:matrix-log}
\index{matrix logarithm}
Let $A \in \Mat_n(\C)$ be invertible (so all eigenvalues are nonzero).
\begin{enumerate}[label=(\alph*)]
  \item For a Jordan block $J_k(\lambda)$ with $\lambda \neq 0$,
    show that $J_k(\lambda) = \lambda(I + \lambda^{-1}N_k)$ and use
    the power series $\log(1 + x) = x - \frac{x^2}{2} +
    \frac{x^3}{3} - \cdots$ (which terminates since $N_k$ is
    nilpotent) to define $\log J_k(\lambda)$.
  \item Compute $\log J_k(\lambda)$ explicitly for $k = 2$ and $k = 3$.
  \item Verify that $e^{\log J_k(\lambda)} = J_k(\lambda)$.
\end{enumerate}
\end{exercise}

\begin{exercise}{Jordan form over $\R$ \challenging}{exo:jordan-real}
\index{Jordan normal form!over $\R$}
The matrix
$A = \begin{pmatrix} 0 & -1 & 1 & 0 \\ 1 & 0 & 0 & 1 \\
0 & 0 & 0 & -1 \\ 0 & 0 & 1 & 0 \end{pmatrix}$
has characteristic polynomial $\chi_A(\lambda) = (\lambda^2+1)^2$.
\begin{enumerate}[label=(\alph*)]
  \item Find the Jordan form of~$A$ over~$\C$.
  \item Show that $A$ is not similar over~$\R$ to a Jordan matrix.
  \item Find the \emph{real Jordan form} of~$A$, i.e.\ a block
    matrix with $2\times 2$ ``rotation-scaling'' blocks
    $\begin{psmallmatrix} a & -b \\ b & a \end{psmallmatrix}$
    replacing complex eigenvalues $a \pm bi$.
\end{enumerate}
\end{exercise}


% ============================================================
\section{Chapter summary}\label{sec:jordan-summary}
% ============================================================

\begin{itemize}
  \item \textbf{Nilpotent endomorphisms.}\index{nilpotent!summary}
    An endomorphism~$f$ is nilpotent if $f^r = 0$ for some $r \geq 1$.
    Its only eigenvalue is~$0$, $\chi_f(\lambda) = (-\lambda)^n$, and
    the index of nilpotency satisfies $r \leq n$.

  \item \textbf{Generalized eigenspaces.}\index{generalized eigenspace!summary}
    For each eigenvalue~$\lambda$ of~$f$, the generalized eigenspace
    $G_\lambda(f) = \Ker\bigl((f - \lambda\Id)^{n_\lambda}\bigr)$ has
    dimension equal to the algebraic multiplicity~$n_\lambda$.  The
    primary decomposition theorem gives
    $E = \bigoplus_{\lambda \in \Spec(f)} G_\lambda(f)$.

  \item \textbf{Jordan blocks.}\index{Jordan block!summary}
    $J_k(\lambda) = \lambda I_k + N_k$ is the basic building block.
    It has eigenvalue~$\lambda$ with algebraic multiplicity~$k$ and
    geometric multiplicity~$1$.

  \item \textbf{Jordan Normal Form.}\index{Jordan normal form!summary}
    Over an algebraically closed field, every endomorphism has a
    unique (up to block ordering) Jordan form.  The block sizes are
    determined by the dimensions $\dim\Ker(f - \lambda\Id)^k$.

  \item \textbf{Minimal polynomial.}
    $\mu_f(\lambda) = \prod_{i} (\lambda - \lambda_i)^{r_i}$,
    where $r_i$ is the size of the \emph{largest} Jordan block for
    eigenvalue~$\lambda_i$.  A matrix is diagonalizable iff
    $\mu_f$ has only simple roots.

  \item \textbf{Matrix exponential.}\index{matrix exponential!summary}
    $e^{tJ_k(\lambda)} = e^{\lambda t} \cdot (\text{upper triangular
    polynomial matrix in~$t$})$.  This gives the explicit solution
    $\vec{x}(t) = e^{tA}\vec{x}_0$ to $\vec{x}' = A\vec{x}$, even
    when $A$ is not diagonalizable.  Non-trivial Jordan blocks produce
    solutions of the form $t^j e^{\lambda t}$ (polynomial times
    exponential).

  \item \textbf{Matrix functions.}
    Any sufficiently smooth function~$h$ can be extended to matrices
    via the Jordan form: $h(J_k(\lambda))$ involves
    $h(\lambda), h'(\lambda), \ldots, h^{(k-1)}(\lambda)$.
\end{itemize}
